%% test-common.tex
%% ===============
%% Helpers shared by the test files; \input this file after
%% \usepackage{showhanmengrid}.
%%
%% \SHGexpect{<dir>}{<chars>}{<lines>}{<along>}{<across>}
%%   Check the grid computed at \begin{document}: <dir> is yoko or tate;
%%   <along> and <across> are the overflow along/across the lines, each
%%   "zero" or "positive" (or "any").  Logs "SHG-EXPECT OK" or
%%   "SHG-EXPECT NG".
%% \SHGmark
%%   Put this at the very beginning of the first line of a page.  When
%%   the page is shipped out, the position of the first character is
%%   compared with the top left square of the grid as drawn; logs
%%   "SHG-POS OK" or "SHG-POS NG".  LuaTeX-ja typesets tate text in a
%%   rotated box whose \savepos coordinates are not those of the page,
%%   so the check is skipped there ("SHG-POS SKIP"); the drawing code is
%%   the same as for pTeX, and test-ptex-tate.tex covers it.
\makeatletter
\ifdefined\savepos
  \let\SHG@savepos\savepos \let\SHG@lastxpos\lastxpos \let\SHG@lastypos\lastypos
  \def\SHG@pageht{\pageheight}
\else
  \let\SHG@savepos\pdfsavepos \let\SHG@lastxpos\pdflastxpos \let\SHG@lastypos\pdflastypos
  \def\SHG@pageht{\pdfpageheight}
\fi
\def\SHG@ok#1{\typeout{SHG-#1 OK}}
\def\SHG@ng#1#2{\typeout{SHG-#1 NG: #2}}

\newcommand\SHGexpect[5]{%
  \AddToHook{begindocument/end}{\SHG@expect{#1}{#2}{#3}{#4}{#5}}}
\ExplSyntaxOn
\cs_new_eq:NN \SHG@ifstreq \str_if_eq:nnTF
\ExplSyntaxOff
\def\SHG@sign#1#2{%
  \SHG@ifstreq{#1}{any}{any}%
    {\ifdim#2>\z@ positive\else\ifdim#2=\z@ zero\else negative\fi\fi}}
\def\SHG@expect#1#2#3#4#5{%
  \edef\SHG@got{\ifshg@tate tate\else yoko\fi/\the\shg@chars/\the\shg@lines/%
    \SHG@sign{#4}\shg@overwd/\SHG@sign{#5}\shg@overht}%
  \edef\SHG@want{#1/#2/#3/#4/#5}%
  \typeout{SHG-RESULT \SHG@got\space(overflow \the\shg@overwd, \the\shg@overht)}%
  \ifx\SHG@got\SHG@want\SHG@ok{EXPECT}\else\SHG@ng{EXPECT}{want \SHG@want}\fi}

\newif\ifSHG@marked
\newcommand\SHGmark{\SHG@savepos\global\SHG@markedtrue}
\def\SHG@near#1#2#3#4{% |#1-#2| < 0.01pt ? #3 : #4
  \ifdim\dimexpr#1-#2\relax<0.01pt
    \ifdim\dimexpr#2-#1\relax<0.01pt #3\else#4\fi
  \else#4\fi}
\AddToHook{shipout/after}{%
  \ifSHG@marked
    \global\SHG@markedfalse
    \ifnum0\ifshg@tate\ifx\shg@engine\SHG@luatexja1\fi\fi=1
      \typeout{SHG-POS SKIP (page \thepage)}%
    \else
      \begingroup
        \dimen0=\SHG@lastxpos sp
        \dimen2=\dimexpr\SHG@pageht-\SHG@lastypos sp\relax
        %% the first character: yoko, its baseline on the left edge;
        %% tate, its baseline at the top edge
        \ifshg@tate
          \dimen4=\dimexpr\shg@gridleft+\shg@gridwd-.5\shg@zw\relax
          \dimen6=\shg@gridtop
        \else
          \dimen4=\shg@gridleft
          \dimen6=\dimexpr\shg@gridtop+\shg@ascent\relax
        \fi
        \SHG@near{\dimen0}{\dimen4}{\SHG@near{\dimen2}{\dimen6}%
          {\SHG@ok{POS}}{\SHG@ng{POS}{page \thepage: y=\the\dimen2, want \the\dimen6}}}%
          {\SHG@ng{POS}{page \thepage: x=\the\dimen0, want \the\dimen4}}%
      \endgroup
    \fi
  \fi}
\def\SHG@luatexja{luatexja}

%% Body text
\newcommand\SHGtext{%
  本は、文字を並べた紙を束ねたものです。
  文字をどこに、どの大きさで、どれだけの間隔で並べるかによって、同じ文章でも読みやすさはまったく変わります。
  版面は、文字の大きさと、1行に並べる文字の数と、1ページに並べる行の数と、行と行の間隔から組み立てます。
  この4つが決まれば、版面の幅と高さは計算で決まります。\par}
\makeatother
\endinput
